We demonstrated that our proposed iterative method effectively reconstructs population trajectories from sample snapshots at multiple time points, achieving greater accuracy and lower computational cost compared to existing approaches. A key advantage of our method is the reuse of a single refined reference across all intervals, allowing temporal information to propagate from one drift update to the next. However, when the data contain too few time points or exhibit weak temporal structure, one should not expect substantial improvements over pairwise methods --- though our approach should not perform worse in these cases.

Several interesting directions remain for future work. First, determining the precise conditions under which a latent SDE is identifiable from marginal samples remains an open problem; see \cref{sec:identifiability} for further discussion. Second, while our current method treats the volatility $\volatility$ as fixed and constant across each iteration, one could introduce an outer loop that iteratively re-estimates $\volatility$ from the data. In this approach, each iteration of our algorithm would still use a fixed $\volatility$, preserving the standard KL-based formulation. Afterward, one would update $\volatility$ --- for instance, by matching the residual variance of the SB-imputed trajectories --- before the next round of inference. However, it is unclear how this extra loop would affect overall convergence, since each update to $\volatility$ changes the underlying SB optimization problem. This is an interesting direction for future work.

Additionally, our iterative refinement procedure may converge to suboptimal local minima. In practice, incorporating multiple random initializations or performing sensitivity analyses could help assess the robustness of the solution, and we plan to explore these aspects in future work. Our method could also be combined with existing approaches (including \dmsbname{} or TrajectoryNet) to enforce constraints such as trajectory smoothness. While our current formulation assumes no observational noise, we may incorporate noise using ideas from \citet{Wang2023}. Finally, we hope to extend our method to make forecasts, allowing predictions beyond the observed time points.